\clearpage
\fancyhead[L]{\small ANZEIGEN.}
\hypertarget{native-512-lower}{}\label{native-512-lower}
\begin{center}\small Göttingische gelehrte Anzeigen. Stück 164. S. 1639..1640. 13. Oktober 1808.\\[3pt]\rule{.72\linewidth}{.3pt}\end{center}

Genf. In der Bibliothèque Britannique vom Mai 1808 finden wir unter der Überschrift \emph{Sur les vingt-une dernières Comètes et les nouvelles planètes; par le Prof.} \textsc{P. Picot} einen Auszug aus einem Brief von \textsc{Olbers} mit einigen (bei uns hinlänglich bekannten) Nachrichten über den jüngsten Kometen und die vier neuen Planeten. Wertvoll ist die Zusammenstellung der parabolischen Elemente der 21 jüngsten seit 1790 erschienenen Kometen nach den zuverlässigsten Bestimmungen. \emph{Neun} davon stammen von \textsc{Olbers} selbst. Zugleich lässt uns dieser vortreffliche Astronom auf eine eigene Schrift über seine bekannte Hypothese zur Entstehung der Asteroiden hoffen. --- Ferner \emph{Demonstration de la création immédiate de la terre en état solide, et de l'impossibilité des causes physiques pour la formation de sa figure, par Mr. l'Abbé} \textsc{Sigorgne}. Da sowohl der von \textsc{Lacaille} am Kap der Guten Hoffnung gemessene Breitengrad als auch die verschiedenen auf der Nordhalbkugel gemessenen Breitengrade eine unregelmäßige Gestalt der Erde bewiesen hätten, könne die Erde weder flüssig noch weich gewesen sein, als sie ihre Drehbewegung erhielt. Sie müsse damals schon ein fester Körper und folglich auch schon vorher an den Polen abgeplattet gewesen sein. Der Rezensent hält weder die Prämissen für gesichert noch die Schlussfolgerung für richtig. \textsc{Lacaille}s Breitengrad dürfte ebenso sehr einer neuen Korrektur bedürfen wie der lappländische damals. Selbst der in Peru gemessene ist noch mit Zweifeln behaftet. Zwar lässt sich nicht bezweifeln, dass der Erdkörper bis zu einem gewissen Grad unregelmäßig ist. Daraus kann man aber keineswegs schließen, dass eine Entstehung durch physikalische Ursachen unmöglich sei. Die Erde hätte nur dann ein vollkommen regelmäßiger Körper werden können, wenn alle Teile der Masse, aus der sie sich bildete, vollkommen flüssig gewesen und so lange geblieben wären, bis sie sich nach den Gesetzen der Hydrostatik angeordnet hätten.
